\subsection{Core Functions}

The listings below show the core functions realising \eqref{eq:fspl}--\eqref{eq:sinr} and the scintillation fade-depth models.

%\begin{figure}[!htbp]
\begin{lstlisting}[caption={Core channel-model functions},label={lst:core}]
import numpy as np
C = 3e8  # speed of light [m/s]

# (1) HAPS free-space path loss  ->  eq. (1)
def haps_fspl_db(f_mhz, d_km):
    return 32.44 + 20*np.log10(f_mhz) + 20*np.log10(d_km)

# (1b) HAPS-to-user/gateway air-to-ground path loss -> eq. (1), Arani et al. Eq. 2
# Pure free-space path loss: NO landscape/LoS-NLoS term on the HAPS tier
# (service link 2 GHz AND feeder link 38 GHz both use this -- see S:IV.C).
def haps_a2g_pathloss_db(r_horiz, h_haps, f_hz, h_user=1.5):
    dh  = h_haps - h_user
    d3d = np.sqrt(r_horiz**2 + dh**2)
    return 20*np.log10(4*np.pi*f_hz*d3d/C)

# (2) UAV LoS probability: exact building-blockage series -> eq. (2)-(3)
# UAV/relay tier ONLY (300 m altitude) -- building blockage does not apply
# to the HAPS tier above.
def los_probability(r, h_uav, h_user=1.5, alpha=0.3, beta=500.0, xi=15.0):
    dh = h_uav - h_user
    J  = max(int(np.floor((r/1000.0)*np.sqrt(alpha*beta) - 1.0)), -1)
    P  = 1.0
    for n in range(J + 1):
        frac = (n + 0.5) / (J + 1)
        val  = h_uav - frac*dh          # h_u - frac*(h_u-h_k)
        P   *= 1.0 - np.exp(-(val**2) / (2*xi**2))
    return P

# (3) UAV air-to-ground average path loss  ->  eq. (4)-(5)
# UAV/relay tier ONLY -- LoS/NLoS blend via Holis & Pechac (2008).
def uav_a2g_pathloss_db(r, h_uav, f_hz=2e9, h_user=1.5,
                        eta_los=1.0, eta_nlos=20.0, **env):
    dh   = h_uav - h_user
    d3d  = np.sqrt(r**2 + dh**2)
    plos = los_probability(r, h_uav, h_user, **env)
    base = 20*np.log10(4*np.pi*f_hz*d3d/C)
    return plos*(base+eta_los) + (1-plos)*(base+eta_nlos)

# (4) Received power and noise-limited SINR  ->  eq. (6)
def sinr_db(p_tx_dbm, pl_db, noise_dbm=-100.0):
    p_rx = 10**((p_tx_dbm - pl_db)/10.0)
    return 10*np.log10(p_rx / 10**(noise_dbm/10.0))
\end{lstlisting}
%\end{figure}

%\begin{figure}[!htbp]
\begin{lstlisting}[caption={Scintillation fading model (ITU P.618)},label={lst:scint}]
# Scintillation fade depth (deterministic) -> Equation (12) ITU P.618
def scintillation_fade_depth_db(p_time_percent, elevation_deg, freq_ghz=20.0):
    """ITU P.618 fade depth: D = D_ref * (f/f_ref)^0.7
                                 * (1/sin(elev)) * sqrt(log(1/p))"""
    elev = np.clip(elevation_deg, 5.0, 90.0)
    freq_factor = (freq_ghz / 20.0) ** 0.7
    elev_rad = np.radians(elev)
    elev_factor = 1.0 / np.sin(elev_rad)
    p_frac = p_time_percent / 100.0
    time_factor = np.sqrt(np.abs(np.log(p_frac + 1e-6)))
    ref_fade_db = 2.5  # at p=1%, elev=30deg, f=20GHz
    return ref_fade_db * freq_factor * elev_factor * time_factor

# Random fading envelope (log-normal or Gamma)
def scintillation_fading_envelope(time_samples, p_time_percent, elevation_deg,
                                   freq_ghz=20.0, dist_km=50.0,
                                   distribution='log_normal'):
    """Generate fading multipliers (0-1 linear) matching ITU statistics."""
    fade_depth_db = scintillation_fade_depth_db(p_time_percent, elevation_deg,
                                                freq_ghz)
    fade_depth_linear = 10.0 ** (-fade_depth_db / 20.0)
    p_in_fade = p_time_percent / 100.0
    path_variance_scale = np.sqrt(dist_km / 50.0)

    if distribution == 'log_normal':
        mu = 0.0
        sigma = np.sqrt(2.0 * np.log(1.0 / fade_depth_linear)) * \
                path_variance_scale
        fading_db = np.random.normal(mu, sigma, time_samples)
        fading_linear = 10.0 ** (fading_db / 20.0)
    else:  # Gamma
        shape = 2.0 * path_variance_scale
        scale = 1.0 / (shape * (1.0 - p_in_fade))
        fading_linear = np.random.gamma(shape, scale, time_samples)

    return np.clip(fading_linear, fade_depth_linear, 1.0)

# Total feeder link loss including scintillation
def feeder_link_loss_with_scintillation(distance_km, elevation_deg,
    freq_ghz=20.0, p_time_percent=1.0, include_scint=True):
    """L_total = L_FSPL + D(p,theta,f) -> eq. (14)"""
    fspl_db_val = 20*np.log10(4*np.pi*freq_ghz*1e9*distance_km*1000 / 3e8)
    if include_scint:
        fade_depth_db = scintillation_fade_depth_db(p_time_percent,
                                                    elevation_deg, freq_ghz)
        return fspl_db_val + fade_depth_db
    else:
        return fspl_db_val
\end{lstlisting}
%\end{figure}

%\begin{figure}[!htbp]
\begin{lstlisting}[language=Python,label={lst:atm_api},caption={Atmospheric losses API (atmospheric\_losses.py)}]
# Import functions
from atmospheric_losses import (
    rain_attenuation_db,        # ITU-R P.838 rain model
    fog_attenuation_db,         # ITU-R P.840 fog model
    gas_attenuation_db,         # ITU-R P.676 gas absorption
    scintillation_fading_db,    # ITU-R P.618 scintillation
    total_feeder_loss_db,       # Combined sum of all effects
    get_rainfall_rate_for_availability,  # P.837 climatology
    get_standard_atmosphere,    # Reference atmosphere
)

# Example: Link budget at 38 GHz, 11.3$^\circ$ elevation
losses = total_feeder_loss_db(
    rain_rate_mmhr=10.0,        # Design rain rate [mm/h]
    liquid_water_g_m3=0.05,     # Fog density [g/m^3]
    pressure_hpa=1013.0,        # Barometric pressure
    temp_c=15.0,                # Temperature
    rh_percent=60.0,            # Relative humidity
    freq_ghz=38.0,              # Frequency
    elevation_deg=45.0,         # Elevation angle
    time_percentage=0.01,       # Outage allowance [%]
    lat=28.5, lon=77.2          # Location (Delhi)
)

print(f"Total atmospheric loss: {losses['total']:.2f} dB")
print(f"  Rain:     {losses['rain']:.2f} dB")
print(f"  Fog:      {losses['fog']:.2f} dB")
print(f"  Gas:      {losses['gas']:.2f} dB")
print(f"  Scint:    {losses['scintillation']:.2f} dB")
\end{lstlisting}
%\end{figure}

%\begin{figure}[!htbp]
\begin{lstlisting}[language=Python,label={lst:climate_api},caption={Climate data API}]
from atmospheric_losses import (
    get_rainfall_rate_for_availability,
    get_rainfall_probability,
)

# Get rain rate for 99.99% link availability (0.01% outage)
# This is the "1-in-10000" rain event from 30-year climatology
rain_99_99pct = get_rainfall_rate_for_availability(
    lat=28.5, lon=77.2,
    availability_percent=99.99
)

# Get probability of rain occurrence
prob_rain = get_rainfall_probability(lat=28.5, lon=77.2)

print(f"Design rain rate: {rain_99_99pct:.1f} mm/h (0.01% exceedance)")
print(f"Rain probability: {prob_rain:.1f}% of average year")
\end{lstlisting}
%\end{figure}
